\documentclass[11pt]{article}

\usepackage[a4paper,margin=1in]{geometry}
\usepackage{amsmath,amssymb,amsthm,mathtools}
\usepackage{microtype}

\newtheorem{theorem}{Theorem}
\newtheorem{lemma}[theorem]{Lemma}
\newtheorem{proposition}[theorem]{Proposition}
\newtheorem{remark}[theorem]{Remark}

\newcommand{\dd}{\,\mathrm{d}}
\newcommand{\one}{\mathbf 1}
\newcommand{\ip}[2]{\left\langle #1,#2\right\rangle}

\title{A Goodman-Type Lower Bound for the House Graph\\
in the High-Density Regime}
\author{}
\date{}

\begin{document}

\maketitle

\section{Statement}

Let \(H\) be the house graph: it is obtained by gluing a triangle and a
\(4\)-cycle along an edge.  Equivalently, we may take
\[
V(H)=\{x,y,a,b,z\}
\]
and
\[
E(H)
=
\{xy,xa,ay,xb,bz,zy\}.
\]
Thus \(xya\) is a triangle, while \(xy z b x\) is a \(4\)-cycle, and the
two subgraphs share the edge \(xy\).

For a graphon \(W:[0,1]^2\to[0,1]\), write
\[
p=t(K_2,W)
=
\int_{[0,1]^2}W(x,y)\,\dd x\,\dd y
\]
for its edge density.

The chromatic polynomial of \(H\) is
\[
\chi_H(k)
=
k(k-1)(k-2)(k^2-3k+3).
\]
Indeed, \(H\) is obtained by gluing \(K_3\) and \(C_4\) along \(K_2\), so
\[
\chi_H(k)
=
\frac{\chi_{K_3}(k)\chi_{C_4}(k)}{\chi_{K_2}(k)},
\]
and
\[
\chi_{K_3}(k)=k(k-1)(k-2),
\qquad
\chi_{C_4}(k)
=
(k-1)^4+(k-1)
=
k(k-1)(k^2-3k+3).
\]
Hence, putting
\[
q:=1-p,
\qquad
r:=2p-1,
\qquad
c:=p^3+q^3,
\]
we obtain
\[
q^5\chi_H\!\left(\frac1q\right)
=
p(2p-1)(p^3+q^3)
=
prc.
\]

We shall prove the following.

\begin{theorem}\label{thm:main}
Let \(H\) be the house graph and let \(W\) be a graphon of edge density
\(p=t(K_2,W)\).  If
\[
p\ge \frac{1+1/\sqrt3}{2},
\]
then
\[
t(H,W)
\ge
(1-p)^5
\chi_H\!\left(\frac1{1-p}\right).
\]
Equivalently,
\[
\boxed{
t(H,W)
\ge
p(2p-1)
\bigl(p^3+(1-p)^3\bigr).
}
\]
\end{theorem}

The proof consists of three elementary ingredients:

\begin{enumerate}
    \item Goodman's lower bound for the triangle density;
    \item the sharp Goodman-type lower bound for \(C_5\);
    \item a gluing inequality for three triangles arranged as a clique tree.
\end{enumerate}

We first establish these ingredients.

\section{Goodman's triangle bound}

\begin{lemma}[Goodman]\label{lem:goodman}
For every graphon \(W\) of edge density \(p\),
\[
t(K_3,W)\ge 2p^2-p=p(2p-1).
\]
\end{lemma}

\begin{proof}
Let
\[
d_W(x):=\int_0^1W(x,z)\,\dd z
\]
be the degree function of \(W\).

For every \(a,b\in[0,1]\),
\[
ab\ge a+b-1.
\]
Consequently, for almost every \(x,y\),
\[
\begin{aligned}
\int_0^1 W(x,z)W(y,z)\,\dd z
&\ge
\int_0^1
\bigl(W(x,z)+W(y,z)-1\bigr)\,\dd z\\
&=
d_W(x)+d_W(y)-1.
\end{aligned}
\]
Multiplying by \(W(x,y)\ge0\) and integrating gives
\[
\begin{aligned}
t(K_3,W)
&\ge
\int_{[0,1]^2}
W(x,y)
\bigl(d_W(x)+d_W(y)-1\bigr)
\,\dd x\,\dd y\\
&=
2\int_0^1d_W(x)^2\,\dd x-p.
\end{aligned}
\]
Since
\[
\int_0^1d_W(x)\,\dd x=p,
\]
Cauchy--Schwarz yields
\[
\int_0^1d_W(x)^2\,\dd x\ge p^2.
\]
Therefore
\[
t(K_3,W)\ge 2p^2-p=p(2p-1).
\]
\end{proof}

\section{The sharp \(C_5\) bound}

We next record a self-contained proof of the sharp odd-cycle bound for
\(C_5\).

\begin{lemma}[Sharp \(C_5\) bound]\label{lem:c5}
For every graphon \(W\) of edge density \(p\),
\[
t(C_5,W)
\ge
p^5-p(1-p)^4.
\]
\end{lemma}

\begin{proof}
Put
\[
U:=1-W,
\qquad
q:=t(K_2,U)=1-p.
\]
If \(q\ge1/2\), equivalently \(p\le1/2\), then
\[
p^5-pq^4
=
p(p^4-q^4)
\le0,
\]
and the assertion follows from \(t(C_5,W)\ge0\).  Hence it is enough to
consider
\[
0\le q\le\frac12.
\]

For \(j\ge1\), let
\[
x_j:=t(P_j,U),
\]
where \(P_j\) denotes the path with \(j\) edges, and put
\[
c_5:=t(C_5,U).
\]

Expanding
\[
t(C_5,W)=t(C_5,1-U)
\]
by inclusion--exclusion over the five edges of \(C_5\) gives
\[
t(C_5,1-U)
=
1-5q+5x_2+5q^2
-5x_3-5qx_2+5x_4-c_5.
\]
Since \(0\le U\le1\), deleting one edge from the integrand defining
\(c_5\) gives
\[
c_5\le x_4.
\]
Therefore
\[
t(C_5,1-U)
\ge
1-5q+5q^2
+5(1-q)x_2
-5x_3
+4x_4.
\]
On the other hand,
\[
(1-q)^5-(1-q)q^4
=
1-5q+10q^2-10q^3+4q^4.
\]
Thus it is enough to prove
\begin{equation}\label{eq:c5certificate}
4x_4-5x_3+5(1-q)x_2
-
\bigl(4q^4-10q^3+5q^2\bigr)
\ge0.
\end{equation}

Let \(T=T_U\) be the self-adjoint graphon operator
\[
Tf(x)
=
\int_0^1U(x,y)f(y)\,\dd y.
\]
Write
\[
T\one=q\one+g,
\qquad
\ip{g}{\one}=0.
\]
We compute the first three path moments.

First,
\[
x_2
=
\|T\one\|_2^2
=
q^2+\|g\|_2^2.
\]
Next,
\[
T^2\one
=
T(q\one+g)
=
q^2\one+qg+Tg,
\]
and hence
\[
\begin{aligned}
x_3
&=
\ip{T\one}{T^2\one}\\
&=
\ip{q\one+g}{q^2\one+qg+Tg}\\
&=
q^3+2q\|g\|_2^2+\ip{g}{Tg},
\end{aligned}
\]
where we used
\[
\ip{\one}{Tg}
=
\ip{T\one}{g}
=
\|g\|_2^2.
\]
Finally,
\[
x_4
=
\|T^2\one\|_2^2,
\]
so
\[
x_4
=
q^4
+3q^2\|g\|_2^2
+2q\ip{g}{Tg}
+\|Tg\|_2^2.
\]

Substitution into the left-hand side of
\eqref{eq:c5certificate} gives
\[
4\|Tg\|_2^2
+
(8q-5)\ip{g}{Tg}
+
(12q^2-15q+5)\|g\|_2^2.
\]
Completing the square,
\[
\begin{aligned}
&4\|Tg\|_2^2
+
(8q-5)\ip{g}{Tg}
+
(12q^2-15q+5)\|g\|_2^2\\
&\qquad=
4\left\|
Tg+\frac{8q-5}{8}g
\right\|_2^2
+
\left(
8q^2-10q+\frac{55}{16}
\right)
\|g\|_2^2.
\end{aligned}
\]
The remaining scalar coefficient is strictly positive:
\[
8q^2-10q+\frac{55}{16}
=
8\left(q-\frac58\right)^2+\frac5{16}>0.
\]
Thus \eqref{eq:c5certificate} holds, completing the proof.
\end{proof}

\section{A three-triangle gluing inequality}

We now prove the special clique-tree inequality needed in the main
argument directly in graphon space.

For a graphon \(W\), define its two-step kernel
\[
C_W(x,y)
:=
\int_0^1W(x,s)W(s,y)\,\dd s.
\]

Let \(F\) be the graph on five vertices \(x,y,z,a,b\) whose three maximal
cliques are
\[
\{x,y,z\},\qquad
\{x,y,a\},\qquad
\{x,z,b\}.
\]
Thus \(F\) consists of three triangles, with the first triangle sharing
the edge \(xy\) with the second and the edge \(xz\) with the third.

\begin{lemma}[Three-triangle gluing]\label{lem:triangle-tree}
Let
\[
p=t(K_2,W),
\qquad
\tau=t(K_3,W).
\]
If \(p>0\), then
\[
\boxed{
t(F,W)\ge\frac{\tau^3}{p^2}.
}
\]
\end{lemma}

\begin{proof}
If \(\tau=0\), the claim is immediate.  Assume henceforth that
\(\tau>0\).

Write simply
\[
C(x,y):=C_W(x,y).
\]
Integrating first over the two vertices \(a,b\) gives
\begin{equation}\label{eq:F-density}
t(F,W)
=
\int_{[0,1]^3}
W(x,y)W(x,z)W(y,z)
C(x,y)C(x,z)
\,\dd x\,\dd y\,\dd z.
\end{equation}

Consider the probability density
\[
\rho(x,y,z)
:=
\frac{
W(x,y)W(x,z)W(y,z)
}{\tau}
\]
on \([0,1]^3\).
Its \((x,y)\)-marginal is the probability density
\[
\nu(x,y)
=
\frac{W(x,y)C(x,y)}{\tau}.
\]
The \((x,z)\)-marginal is the same, after relabelling variables.

Also consider the probability density
\[
\mu(x,y)
=
\frac{W(x,y)}{p}
\]
on \([0,1]^2\).
We have \(\nu\ll\mu\), with Radon--Nikodym derivative
\[
\frac{\dd\nu}{\dd\mu}(x,y)
=
\frac{p\,C(x,y)}{\tau}.
\]
Nonnegativity of relative entropy gives
\[
0
\le
D(\nu\Vert\mu)
=
\int_{[0,1]^2}
\nu(x,y)
\log\left(
\frac{pC(x,y)}{\tau}
\right)
\,\dd x\,\dd y.
\]
Therefore
\begin{equation}\label{eq:geom-codegree}
\int\nu(x,y)\log C(x,y)\,\dd x\,\dd y
\ge
\log\frac{\tau}{p}.
\end{equation}

By \eqref{eq:F-density},
\[
\frac{t(F,W)}{\tau}
=
\mathbb E_\rho
\bigl[
C(x,y)C(x,z)
\bigr].
\]
Jensen's inequality for the concave function \(\log\) gives
\[
\begin{aligned}
\log\frac{t(F,W)}{\tau}
&\ge
\mathbb E_\rho
\bigl[
\log C(x,y)+\log C(x,z)
\bigr]\\
&=
2\int\nu(x,y)\log C(x,y)\,\dd x\,\dd y\\
&\ge
2\log\frac{\tau}{p},
\end{aligned}
\]
where the last inequality is \eqref{eq:geom-codegree}.
Exponentiating,
\[
\frac{t(F,W)}{\tau}
\ge
\left(\frac{\tau}{p}\right)^2,
\]
hence
\[
t(F,W)\ge\frac{\tau^3}{p^2}.
\]

All logarithmic expressions above may be interpreted in the usual
extended sense.  Alternatively, one may replace \(C\) by
\(C+\varepsilon\), apply the same argument, and let
\(\varepsilon\downarrow0\).
\end{proof}

\section{The exact two-chord decomposition}

We now turn to the house graph.

Let
\[
U:=1-W.
\]
Again write
\[
C(x,y)
=
C_W(x,y)
=
\int_0^1W(x,s)W(s,y)\,\dd s.
\]
Define
\[
A(x,y):=W(x,y)C(x,y),
\qquad
D(x,y):=U(x,y)C(x,y).
\]
Since \(W+U=1\),
\begin{equation}\label{eq:AplusD}
A(x,y)+D(x,y)=C(x,y).
\end{equation}

Consider the \(5\)-cycle
\[
x-a-y-z-b-x.
\]
After integrating over \(a\) and \(b\), its density is
\[
t(C_5,W)
=
\int_{[0,1]^3}
W(y,z)C(x,y)C(x,z)
\,\dd x\,\dd y\,\dd z.
\]
The two distance-two chords incident with \(x\) are \(xy\) and \(xz\).

Define
\[
\begin{aligned}
F
&:=
\int
W(y,z)A(x,y)A(x,z)
\,\dd x\,\dd y\,\dd z,\\
R
&:=
\int
W(y,z)A(x,y)D(x,z)
\,\dd x\,\dd y\,\dd z,\\
G
&:=
\int
W(y,z)D(x,y)D(x,z)
\,\dd x\,\dd y\,\dd z.
\end{aligned}
\]
The integral with \(D(x,y)A(x,z)\) instead of
\(A(x,y)D(x,z)\) also equals \(R\), by interchanging \(y\) and \(z\).

Expanding \eqref{eq:AplusD} therefore gives
\begin{equation}\label{eq:c5split}
t(C_5,W)=F+2R+G.
\end{equation}

On the other hand, adding the chord \(xy\) to the \(5\)-cycle gives
exactly the house graph: \(xay\) is the triangle and
\(xy z b x\) is the \(4\)-cycle.
Consequently
\[
\begin{aligned}
t(H,W)
&=
\int
W(y,z)A(x,y)C(x,z)
\,\dd x\,\dd y\,\dd z\\
&=
F+R.
\end{aligned}
\]
Combining this with \eqref{eq:c5split}, we obtain the exact identity
\begin{equation}\label{eq:key-decomposition}
\boxed{
2t(H,W)
=
F+t(C_5,W)-G.
}
\end{equation}

This identity is the main structural point of the proof.

\section{Triangle surplus and the bad two-chord term}

Let
\[
\tau:=t(K_3,W)
\]
and define the Goodman surplus
\[
\delta
:=
\tau-p(2p-1)
=
\tau-pr.
\]
By Lemma~\ref{lem:goodman},
\begin{equation}\label{eq:delta-nonneg}
\delta\ge0.
\end{equation}

We shall show that the negative term \(G\) in
\eqref{eq:key-decomposition} is bounded by this same surplus.

Let
\[
d_U(x):=\int_0^1U(x,y)\,\dd y
\]
and
\[
x_2:=t(P_2,U)
=
\int_0^1d_U(x)^2\,\dd x.
\]
Since
\[
\int_0^1d_U(x)\,\dd x=q,
\]
put
\[
v
:=
x_2-q^2
=
\int_0^1(d_U(x)-q)^2\,\dd x
\ge0.
\]
Also define
\begin{equation}\label{eq:n-def}
n
:=
\int_{[0,1]^3}
W(y,z)U(x,y)U(x,z)
\,\dd x\,\dd y\,\dd z.
\end{equation}
Clearly
\[
n\ge0.
\]

Let
\[
\kappa_3:=t(K_3,U).
\]
Splitting the missing edge \(yz\) according to
\(1=W(y,z)+U(y,z)\), we have
\[
x_2=n+\kappa_3,
\]
and hence
\begin{equation}\label{eq:n-x2}
n=x_2-\kappa_3.
\end{equation}

Expanding the triangle density of \(W=1-U\),
\[
\begin{aligned}
\tau
&=
t(K_3,1-U)\\
&=
1-3q+3x_2-\kappa_3.
\end{aligned}
\]
Since
\[
pr
=
(1-q)(1-2q)
=
1-3q+2q^2,
\]
we obtain
\[
\begin{aligned}
\delta
&=
\tau-pr\\
&=
3x_2-\kappa_3-2q^2\\
&=
2(x_2-q^2)+(x_2-\kappa_3).
\end{aligned}
\]
Using the definitions of \(v\) and \(n\),
\begin{equation}\label{eq:delta-decomp}
\boxed{
\delta=2v+n.
}
\end{equation}
In particular,
\begin{equation}\label{eq:n-le-delta}
0\le n\le\delta.
\end{equation}

Now \(0\le C(x,y)\le1\), and therefore
\[
0\le D(x,y)
=
U(x,y)C(x,y)
\le U(x,y).
\]
It follows directly from the definition of \(G\) that
\[
\begin{aligned}
G
&=
\int
W(y,z)
D(x,y)D(x,z)
\,\dd x\,\dd y\,\dd z\\
&\le
\int
W(y,z)
U(x,y)U(x,z)
\,\dd x\,\dd y\,\dd z\\
&=
n.
\end{aligned}
\]
Together with \eqref{eq:n-le-delta}, this gives
\begin{equation}\label{eq:G-delta}
\boxed{
G\le n\le\delta.
}
\end{equation}

\section{Completion of the proof}

We can now prove Theorem~\ref{thm:main}.

\begin{proof}[Proof of Theorem~\ref{thm:main}]
Assume
\[
p\ge\frac{1+1/\sqrt3}{2}.
\]
In particular \(p>1/2\), so
\[
r:=2p-1>0.
\]
Let
\[
q:=1-p,
\qquad
c:=p^3+q^3,
\qquad
B(p):=prc.
\]
Thus \(B(p)\) is precisely the claimed chromatic-polynomial lower bound.

Let
\[
\tau=t(K_3,W)=pr+\delta,
\qquad
\delta\ge0.
\]

The graph counted by \(F\) is exactly the graph of
Lemma~\ref{lem:triangle-tree}, so
\[
F\ge\frac{\tau^3}{p^2}.
\]
Therefore
\[
\begin{aligned}
F-pr^3
&\ge
\frac{(pr+\delta)^3}{p^2}-pr^3\\
&=
3r^2\delta
+
\frac{3r}{p}\delta^2
+
\frac1{p^2}\delta^3.
\end{aligned}
\]
Thus
\begin{equation}\label{eq:F-surplus}
F-pr^3
\ge
3r^2\delta
+
\frac{3r}{p}\delta^2
+
\frac1{p^2}\delta^3.
\end{equation}

By Lemma~\ref{lem:c5},
\begin{equation}\label{eq:C5-surplus}
t(C_5,W)
\ge
p^5-pq^4.
\end{equation}
Since \(p+q=1\) and \(r=p-q\),
\[
p^5-pq^4
=
p(p^4-q^4)
=
pr(p^2+q^2).
\]
Moreover,
\[
\begin{aligned}
r^2+p^2+q^2
&=
(p-q)^2+p^2+q^2\\
&=
2(p^2-pq+q^2)\\
&=
2(p^3+q^3),
\end{aligned}
\]
where in the last equality we used
\[
p^3+q^3=(p+q)^3-3pq(p+q)=1-3pq
=p^2-pq+q^2.
\]
Consequently,
\begin{equation}\label{eq:target-split}
pr^3+\bigl(p^5-pq^4\bigr)
=
pr\bigl(r^2+p^2+q^2\bigr)
=
2pr(p^3+q^3)
=
2B(p).
\end{equation}

Subtracting \eqref{eq:target-split} from the exact decomposition
\eqref{eq:key-decomposition}, we obtain the exact defect identity
\begin{equation}\label{eq:defect-exact}
\boxed{
\begin{aligned}
2\bigl(t(H,W)-B(p)\bigr)
={}&
\bigl(F-pr^3\bigr)\\
&+
\bigl(t(C_5,W)-(p^5-pq^4)\bigr)
-G.
\end{aligned}
}
\end{equation}

The middle term is nonnegative by
\eqref{eq:C5-surplus}.  Combining
\eqref{eq:F-surplus} and \eqref{eq:G-delta} with
\eqref{eq:defect-exact} gives
\[
\begin{aligned}
2\bigl(t(H,W)-B(p)\bigr)
&\ge
3r^2\delta
+
\frac{3r}{p}\delta^2
+
\frac1{p^2}\delta^3
-G\\
&\ge
3r^2\delta
+
\frac{3r}{p}\delta^2
+
\frac1{p^2}\delta^3
-\delta\\
&=
(3r^2-1)\delta
+
\frac{3r}{p}\delta^2
+
\frac1{p^2}\delta^3.
\end{aligned}
\]
Thus we have proved the quantitative estimate
\begin{equation}\label{eq:quantitative-final}
\boxed{
2\bigl(t(H,W)-B(p)\bigr)
\ge
(3r^2-1)\delta
+
\frac{3r}{p}\delta^2
+
\frac1{p^2}\delta^3.
}
\end{equation}

Finally, the hypothesis
\[
p\ge\frac{1+1/\sqrt3}{2}
\]
is equivalent, because \(r=2p-1\ge0\), to
\[
r\ge\frac1{\sqrt3},
\]
and hence
\[
3r^2-1\ge0.
\]
Every term on the right-hand side of
\eqref{eq:quantitative-final} is therefore nonnegative, since
\[
p>0,\qquad r\ge0,\qquad \delta\ge0.
\]
It follows that
\[
t(H,W)\ge B(p)
=
p(2p-1)(p^3+q^3).
\]
Since \(q=1-p\),
\[
t(H,W)
\ge
p(2p-1)
\bigl(p^3+(1-p)^3\bigr).
\]
As shown at the beginning,
\[
p(2p-1)
\bigl(p^3+(1-p)^3\bigr)
=
(1-p)^5
\chi_H\!\left(\frac1{1-p}\right).
\]
This proves the theorem.
\end{proof}

\section{A useful quantitative form}

The proof actually establishes a little more than the theorem.

\begin{proposition}\label{prop:quantitative}
Let \(W\) be any graphon with \(p>1/2\).  Put
\[
r=2p-1
\]
and
\[
\delta
=
t(K_3,W)-p(2p-1)\ge0.
\]
Then
\[
\boxed{
\begin{aligned}
2\Bigl(
t(H,W)
-
p(2p-1)
\bigl(p^3+(1-p)^3\bigr)
\Bigr)
\ge{}&
(3(2p-1)^2-1)\delta\\
&+
\frac{3(2p-1)}{p}\delta^2
+
\frac1{p^2}\delta^3.
\end{aligned}
}
\]
In particular, the desired house inequality follows whenever
\[
3(2p-1)^2\ge1.
\]
\end{proposition}

\begin{remark}
The argument is sharp at every balanced complete \(k\)-partite graphon
whose edge density lies in the proved range.  At such a graphon,
\[
p=1-\frac1k,
\]
the Goodman triangle surplus satisfies
\[
\delta=0,
\]
the \(C_5\) inequality is an equality, and the term \(G\) is zero.
Thus no positive error is introduced at the Tur\'an equality examples.
For the present density range, this includes every \(k\ge5\).
\end{remark}

\end{document}